\chapter{हॉर्मोन वाला संवाद}\label{ch:g8:hormonal-communication}

\omterm{def:g8:puberty:puberty}{यौवनारंभ} ने इनसे मिलवाया था (\cref{def:g8:puberty:hormones}); गोली इनका
फ़ायदा उठाती है (\cref{def:g8:contraception:pill}); और डर की उस गरम लहर
ने दौड़ के बीच एक चिट्ठी पहुँचाई थी
(\cref{ex:g8:nervous-communication:cases})। अब उस डाक-सेवा का अपना दौरा
बनता है: वे \omterm{def:g8:hormonal-communication:glands}{ग्रंथियाँ} जो चिट्ठियाँ लिखती हैं, वह \omterm{prop:g4:journey-of-food:absorb}{रक्त} जो उन्हें ढोता है,
पता लिखने की वह तरकीब जो एक प्रसारण को भी ठीक-ठीक बना देती है --- और वे
\omterm{prop:g8:hormonal-communication:feedback}{प्रतिपुष्टि} वाले फंदे, जो पूरे पत्राचार को संतुलन में रखते हैं।

\section{ग्रंथियाँ, हॉर्मोन, लक्ष्य}

\begin{definition}[अंतःस्रावी ग्रंथियाँ]\label{def:g8:hormonal-communication:glands}
जो \emph{ग्रंथि}\index{ग्रंथि} अपना उत्पाद किसी नली में उतारने के बजाय
(जैसे \omterm{def:g4:journey-of-food:tract}{लार} वाली उतारती है) \emph{\omterm{prop:g4:journey-of-food:absorb}{रक्त} में} छोड़ती है वह
\emph{अंतःस्रावी} ग्रंथि है, और उसके उत्पाद ही
\emph{हॉर्मोन}\index{हॉर्मोन} हैं। मुख्य लिखने वाली: \omterm{def:g5:brain-in-command:system}{मस्तिष्क} के तल की
\emph{कमान-ग्रंथि} --- यानी \omterm{def:g8:puberty:puberty}{यौवनारंभ} को शुरू करने वाली, और बाक़ी
ज़्यादातर की संचालक; गरदन में \emph{थायरॉइड} --- जिसका हॉर्मोन पूरे शरीर
की चाल तय करता है; \omterm{def:g7:removing-wastes:kidneys}{गुर्दों} पर टोपी की तरह बैठी \emph{अधिवृक्क} ग्रंथियाँ
--- यानी डरों और दौड़ों का चेतावनी-हॉर्मोन; \emph{अग्न्याशय} --- जिसका
हॉर्मोन \omterm{prop:g4:journey-of-food:absorb}{रक्त} की \omterm{prop:g7:organs-at-work:bill}{ग्लूकोज़} सँभालता है; और \emph{\omterm{def:g8:reproductive-systems:female}{अंडाशय} तथा \omterm{def:g8:reproductive-systems:male}{वृषण}} ---
यानी \cref{ch:g8:reproductive-systems} वाले चक्र और दोबारा बनावट के
हॉर्मोन।
\end{definition}

\begin{proposition}[ग्राही से पता]\label{prop:g8:hormonal-communication:address}
\omterm{prop:g4:journey-of-food:absorb}{रक्त} हर \omterm{def:g8:hormonal-communication:glands}{हॉर्मोन} को हर जगह पहुँचाता है (\cref{ex:g8:puberty:sites}) ---
फिर भी हर एक सिर्फ़ अपने \emph{लक्ष्य अंगों}\index{लक्ष्य अंग} पर असर
करता है: यानी उन \omterm{def:g5:first-look-at-cells:cell}{कोशिकाओं} पर, जो उसी संदेशवाहक के लिए जुड़ने की जगहें,
\emph{ग्राही}, लिए बनी हैं। ढुलाई प्रसारण से, और पता पाने वाले की तरफ़
से: यानी वह चिट्ठी, जो हर घर को मिलती है और सिर्फ़ वहीं पढ़ी जाती है
जहाँ ताली बैठती हो। (\cref{def:g8:nervous-communication:synapse} वाली
\omterm{def:g8:nervous-communication:synapse}{तंत्रिका-संधि} का जुड़ना भी वही तरकीब है, नैनोमीटर की दूरी पर।)
\end{proposition}
\admitted

\begin{example}[एक संदेशवाहक का दौरा: चेतावनी-हॉर्मोन]\label{ex:g8:hormonal-communication:alarm}
साइकिल पर बाल-बाल बचना: अधिवृक्क \omterm{def:g8:hormonal-communication:glands}{ग्रंथियाँ} अपना चेतावनी-हॉर्मोन \omterm{prop:g4:journey-of-food:absorb}{रक्त} में
उँडेल देती हैं, और सेकंडों में --- यानी \omterm{prop:g4:journey-of-food:absorb}{रक्त} के एक ही चक्कर में
(\cref{ex:g4:heart-and-blood:round}) --- हर उस पाने वाले का, जिस पर वह
ताली बैठती है, फ़ौरन जवाब आता है: \omterm{def:g4:heart-and-blood:heart}{दिल} की चाल उछलती है, हवा-मार्ग चौड़े
होते हैं, \omterm{prop:g7:digestion-and-nutrients:absorption}{यकृत} अपनी जमा की हुई \omterm{prop:g7:organs-at-work:bill}{ग्लूकोज़} छोड़ता है
(\cref{ex:g7:digestion-and-nutrients:breadtrace} वाला बैंक, तलब किया
हुआ), \omterm{def:g3:bones-and-muscles:muscle}{पेशियों} की \omterm{def:g4:heart-and-blood:vessels}{वाहिकाएँ} चौड़ी होती हैं और पाचन-नली की सिकुड़ती हैं
(\cref{prop:g7:organs-at-work:supply} वाला रुख़ मोड़ना, \omterm{def:g8:hormonal-communication:glands}{हॉर्मोन} से थामा
हुआ), और \omterm{def:g1:the-five-senses:sense}{त्वचा} पीली पड़कर झुनझुनाने लगती है। एक चिट्ठी, छह जवाब, और एक
ही तालमेल वाली हालत: तैयार।
\end{example}

\section{प्रतिपुष्टि से संतुलन}

\begin{proposition}[प्रतिपुष्टि का उसूल]\label{prop:g8:hormonal-communication:feedback}
\omterm{def:g8:hormonal-communication:glands}{हॉर्मोन} के तंत्र ख़ुद को \emph{प्रतिपुष्टि}\index{प्रतिपुष्टि} से सँभालते
हैं: लिखने वाली \omterm{def:g8:hormonal-communication:glands}{ग्रंथि} \emph{अपनी ही चिट्ठियों का नतीजा देखती रहती है},
और नतीजा आ जाने पर धीमी पड़ जाती है --- यानी तापनियंत्रक, नल नहीं। इसकी
सबसे अच्छी मिसाल \omterm{prop:g4:journey-of-food:absorb}{रक्त} की \omterm{prop:g7:organs-at-work:bill}{ग्लूकोज़} है: खाने के बाद चढ़ती हुई \omterm{prop:g7:organs-at-work:bill}{ग्लूकोज़}
ख़ुद अग्न्याशय का जमा करने वाला \omterm{def:g8:hormonal-communication:glands}{हॉर्मोन} चला देती है; \omterm{prop:g7:organs-at-work:bill}{ग्लूकोज़} बैंक में
गई, स्तर गिरा, और उसी के साथ वह स्राव भी थम गया; और खानों के बीच एक
साथी \omterm{def:g8:hormonal-communication:glands}{हॉर्मोन} उसी बैंक को दूसरी तरफ़ से खोलता है। वह स्तर दिन-रात एक तंग
गली में चलता रहता है, और कोई होशमंद चलाने वाला होता ही नहीं।
\end{proposition}
\admitted

\begin{example}[जब प्रतिपुष्टि नाकाम होती है: मधुमेह]\label{ex:g8:hormonal-communication:diabetes}
\emph{मधुमेह}\index{मधुमेह} में जमा करने वाले \omterm{def:g8:hormonal-communication:glands}{हॉर्मोन} का इशारा नाकाम हो
जाता है --- या तो अग्न्याशय उसे लिखना बंद कर देता है, या पाने वाले जवाब
देना। हर खाने के बाद \omterm{prop:g7:organs-at-work:bill}{ग्लूकोज़} \omterm{prop:g4:journey-of-food:absorb}{रक्त} में ढेर होती जाती है; \omterm{def:g7:removing-wastes:kidneys}{गुर्दे} का वापस
लेना दब जाता है और \omterm{def:g7:removing-wastes:kidneys}{मूत्र} में शर्करा आ जाती है --- यानी
\cref{ex:g7:removing-wastes:thrift} वाला चिर-परिचित संकेत, आख़िरकार
समझाया हुआ। और इलाज व्यावहारिक अध्याय है: ग़ायब \omterm{def:g8:hormonal-communication:glands}{हॉर्मोन} की नपी-तुली
ख़ुराकें, खानों के हिसाब से सधी हुईं --- यानी सुई से लिखी हुई चिट्ठी।
लाखों लोग ठीक इसी इंतज़ाम पर अच्छी ज़िंदगी जीते हैं।
\end{example}

\begin{remark}[दोनों तंत्रों की साझेदारी]\label{rem:g8:hormonal-communication:partnership}
\Cref{prop:g8:nervous-communication:compare} वाले तार और डाक एक ही
प्रशासन हैं। \omterm{def:g5:brain-in-command:system}{मस्तिष्क} के तल की कमान-ग्रंथि उनका जोड़ है: ऊपर \omterm{def:g5:brain-in-command:system}{तंत्रिका
तंत्र}, नीचे \omterm{def:g8:hormonal-communication:glands}{हॉर्मोन} --- साइकिल वाले का डर \omterm{def:g8:nervous-communication:neuron}{तंत्रिका} से शुरू हुआ और
\omterm{def:g8:hormonal-communication:glands}{हॉर्मोन} से टिका रहा; \omterm{def:g8:puberty:puberty}{यौवनारंभ} \omterm{def:g5:brain-in-command:system}{मस्तिष्क} वाले संचालक के नीचे \omterm{def:g8:hormonal-communication:glands}{हॉर्मोन} की
धीमी रफ़्तार से चलता है; और
\cref{ex:g8:reproduction-and-environment:lambs} वाली भेड़ें दिन की
लंबाई (यानी \omterm{def:g8:nervous-communication:neuron}{तंत्रिका} वाला बोध) को प्रजनन के हॉर्मोनों में बदल देती हैं।
जहाँ रफ़्तार चाहिए वहाँ तेज़ तार, जहाँ कार्यक्रम महीनों चलना हो वहाँ धीमी
डाक, और बीच में अनुवाद के दफ़्तर।
\end{remark}

\section{किसी हॉर्मोन तंत्र को पढ़ना}

\begin{method}[किसी भी हॉर्मोन के चार सवाल]\label{met:g8:hormonal-communication:read}
अपने सामने आए किसी भी \omterm{def:g8:hormonal-communication:glands}{हॉर्मोन} के लिए यह तय करो:
\begin{enumerate}
  \item \emph{लिखने वाला}: कौन-सी \omterm{def:g8:hormonal-communication:glands}{ग्रंथि}, \omterm{prop:g4:journey-of-food:absorb}{रक्त} में, और किस इशारे
        पर?
  \item \emph{पाने वाले}: किन \omterm{prop:g8:hormonal-communication:address}{लक्ष्य अंगों} पर उसके ग्राही हैं, और
        चिट्ठी पहुँचने पर हर एक क्या करता है?
  \item \emph{\omterm{prop:g8:hormonal-communication:feedback}{प्रतिपुष्टि}}: कौन-सा नतीजा वापस ख़बर करता है, और
        ज़्यादा लिख देने से कैसे बचा जाता है?
  \item \emph{नाकामी और इस्तेमाल}: चिट्ठी नाकाम हो तो क्या होता है
        --- और क्या दवा उसकी नक़ल करती है (वही गोली, वही \omterm{ex:g8:hormonal-communication:diabetes}{मधुमेह} वाली
        ख़ुराकें)?
\end{enumerate}
\end{method}

\begin{example}[यह तरीक़ा थायरॉइड की चाल तय करने वाले पर]\label{ex:g8:hormonal-communication:thyroid}
लिखने वाला: थायरॉइड, और उसका संचालन कमान-ग्रंथि करती है। पाने वाले:
क़रीब-क़रीब हर \omterm{def:g5:first-look-at-cells:cell}{कोशिका} --- क्योंकि वह \omterm{def:g8:hormonal-communication:glands}{हॉर्मोन}
\cref{prop:g7:what-is-respiration:power} वाली धीमी आग की दर तय करता है।
\omterm{prop:g8:hormonal-communication:feedback}{प्रतिपुष्टि}: कमान-ग्रंथि उसका स्तर पढ़ती है और अपना संचालन छाँट लेती है।
नाकामी: बहुत कम हो, तो पूरा इंजन सुस्ताने लगता है --- थकान, ठंड, और हर
चीज़ धीमी; और बहुत ज़्यादा हो, तो वह दौड़ने लगता है। इस्तेमाल: ग़ायब
\omterm{def:g8:hormonal-communication:glands}{हॉर्मोन} रोज़ की एक गोली है --- यानी दवा की सबसे पुरानी डाक-चिट्ठियों में
से एक।
\end{example}

\section{अभ्यास}

\begin{exercise}[$\star$]\label{exo:g8:hormonal-communication:1}
कोई \omterm{def:g8:hormonal-communication:glands}{ग्रंथि} अंतःस्रावी किस बात से बनती है? चार के नाम बताओ, और हर एक के
साथ एक-एक \omterm{def:g8:hormonal-communication:glands}{हॉर्मोन} की भूमिका भी।
\end{exercise}

\begin{exercise}[$\star$]\label{exo:g8:hormonal-communication:2}
\omterm{prop:g4:journey-of-food:absorb}{रक्त} हर जगह पहुँचाता है; और \omterm{def:g8:hormonal-communication:glands}{हॉर्मोन} कहीं-कहीं असर करते हैं। इस उलझन को
सुलझाओ।
\end{exercise}

\begin{exercise}[$\star$]\label{exo:g8:hormonal-communication:3}
चेतावनी-हॉर्मोन के छह जवाबों में से चार गिनाओ, और हर एक के साथ उसकी सेवा
करता अध्याय भी।
\end{exercise}

\begin{exercise}[$\star$]\label{exo:g8:hormonal-communication:4}
\omterm{prop:g8:hormonal-communication:feedback}{प्रतिपुष्टि} का उसूल बताओ, और उसकी सबसे अच्छी मिसाल भी।
\end{exercise}

\begin{exercise}[$\star$]\label{exo:g8:hormonal-communication:5}
\omterm{ex:g8:hormonal-communication:diabetes}{मधुमेह} में क्या नाकाम होता है, और उसकी वजह से \omterm{def:g7:removing-wastes:kidneys}{गुर्दे} पर क्या दोबारा
दिखने लगता है?
\end{exercise}

\begin{exercise}[$\star$]\label{exo:g8:hormonal-communication:6}
\omterm{def:g8:nervous-communication:neuron}{तंत्रिका} और \omterm{def:g8:hormonal-communication:glands}{हॉर्मोन} वाले तंत्र कहाँ जुड़ते हैं? किसी एक रिले की मिसाल
दो।
\end{exercise}

\begin{exercise}[$\star\star$]\label{exo:g8:hormonal-communication:7}
\omterm{prop:g8:hormonal-communication:feedback}{प्रतिपुष्टि} के लिए तापनियंत्रक सही तस्वीर क्यों है, और नल ग़लत क्यों?
\end{exercise}

\begin{exercise}[$\star\star$]\label{exo:g8:hormonal-communication:8}
\cref{met:g8:hormonal-communication:read} को अग्न्याशय के जमा करने वाले
\omterm{def:g8:hormonal-communication:glands}{हॉर्मोन} पर चलाओ।
\end{exercise}

\begin{exercise}[$\star\star$]\label{exo:g8:hormonal-communication:9}
\omterm{ex:g8:hormonal-communication:diabetes}{मधुमेह} वाले की सुई और \omterm{def:g8:contraception:pill}{गर्भनिरोधक गोली} को तरीक़े के चौथे सवाल की मदद से
एक ही तरह के काम की तरह समझाओ।
\end{exercise}

\begin{exercise}[$\star\star$]\label{exo:g8:hormonal-communication:10}
इन्हें कारणों के साथ तंत्र सौंपो: कँपकँपी का शुरू होना; दाढ़ी का उगना;
\omterm{prop:g4:journey-of-food:absorb}{रक्त} की \omterm{prop:g7:organs-at-work:bill}{ग्लूकोज़} वाली तंग गली; और घुटने की उछाल।
\end{exercise}

\begin{exercise}[$\star\star$]\label{exo:g8:hormonal-communication:11}
\cref{ex:g8:reproduction-and-environment:lambs} वाली भेड़ें \omterm{def:g6:exploring-our-environment:conditions}{प्रकाश} को
मेमनों में बदल देती हैं।
\cref{rem:g8:hormonal-communication:partnership} के दफ़्तरों से होकर उस
रिले का पीछा करो।
\end{exercise}

\begin{exercise}[$\star\star\star$]\label{exo:g8:hormonal-communication:12}
``ग्राही किसी प्रसारण को निजी लाइनों के जाल में बदल देते हैं।'' इस वाक्य
का बचाव करो, और फिर उसे आगे बढ़ाओ: किसी \omterm{def:g8:hormonal-communication:glands}{हॉर्मोन} के श्रोताओं में
\emph{शामिल} होने के लिए \omterm{def:g5:first-look-at-cells:cell}{कोशिका} में क्या बदलना चाहिए --- और इससे इस बारे
में क्या इशारा मिलता है कि \omterm{def:g8:puberty:puberty}{यौवनारंभ} के लक्ष्य सुनने वाले बने कैसे?
\end{exercise}

\section{समस्या: ग्लूकोज़ की बही}

\begin{problem}[{सप्ताहांत समस्या --- रक्त की शर्करा का एक दिन,
पत्राचार की तरह पढ़ा हुआ}]\label{pb:g8:hormonal-communication:1}
लगातार नापने वाले किसी ग्लूकोज़-यंत्र का (कल्पित) दिन-लेखा: रात भर
एक-सा; 8 बजे नाश्ता --- एक चढ़ाई, और फिर 10 बजे तक शांत उतार; 12 बजे
खेल --- एक छोटा-सा गड्ढा, जो फ़ौरन सँभाल लिया गया; दोपहर का खाना छूट
गया --- फिर भी पूरी दोपहर एक-सा; 17 बजे एक डर (साइकिल पर बाल-बाल बचना)
--- एक तीखी, छोटी चढ़ाई; और 19 बजे रात का खाना --- फिर चढ़ाई और उतार।

\medskip\noindent\textbf{भाग I --- सुबह।}
\begin{enumerate}
  \item रात भर की एक-सी हालत: मालिक के सोते रहते हुए वह गली किस
        दो-चिट्ठियों वाले पत्राचार ने थामी रखी?
  \item नाश्ते की चढ़ाई और उतार: इशारा, लिखने वाला, चिट्ठी, और वे
        पाने वाले भी बताओ जिन्होंने फ़ालतू बैंक में डाला।
  \item उस उतार का \emph{रुक जाना} --- यानी गली से नीचे कोई गोता न
        लगना --- ख़ुद \omterm{prop:g8:hormonal-communication:feedback}{प्रतिपुष्टि} की पहचान क्यों है?
  \item 10 बजे का स्तर 7 बजे वाले से मेल खाता है। तापनियंत्रक वाली
        तस्वीर बीच के स्राव की चाल के बारे में क्या भविष्यवाणी करती
        है?
\end{enumerate}

\medskip\noindent\textbf{भाग II --- दोपहर।}
\begin{enumerate}[resume]
  \item खेल का गड्ढा और उसका सँभलना: किस चिट्ठी ने \omterm{prop:g7:digestion-and-nutrients:absorption}{यकृत} का बैंक तलब
        किया, और किस अध्याय का खाता ख़र्च हो रहा था?
  \item छूटे हुए दोपहर के खाने से दिखने लायक़ कुछ नहीं बदला। किस साथी
        \omterm{def:g8:hormonal-communication:glands}{हॉर्मोन} के चुपचाप काम ने पूरी दोपहर पुल बाँधे रखा?
  \item डर की तीखी चढ़ाई: लिखने वाले का नाम बताओ, और उस अचानक शर्करा
        का मक़सद भी --- \cref{ex:g8:hormonal-communication:alarm} के
        अनुसार, किस चीज़ के लिए तैयार?
  \item \omterm{ex:g8:hormonal-communication:diabetes}{मधुमेह} वाले टिकाऊ ढेर के उलट, डर की वह चढ़ाई घंटे भर में क्यों
        बैठ गई?
\end{enumerate}

\medskip\noindent\textbf{भाग III --- परामर्श।}
\begin{enumerate}[resume]
  \item \omterm{ex:g8:hormonal-communication:diabetes}{मधुमेह} वाले किसी व्यक्ति की उसी दिन की बही हर खाने पर अलग
        होती। बिना इलाज वाले नाश्ते का वक्र बताओ, और वह जगह भी जहाँ
        \omterm{def:g7:removing-wastes:kidneys}{गुर्दे} का संकेत दिखता।
  \item और इलाज वाला रूप: हर खाने से पहले की वह नपी-तुली ख़ुराक किसकी
        नक़ल करती है, और उसका आकार थाली से क्यों मेल खाना चाहिए?
  \item डॉक्टर उस स्वस्थ बही को ``वह बातचीत, जो तुम कभी सुनते नहीं''
        कहते हैं। उस दिन के सारे बोलने वाले --- \omterm{def:g8:hormonal-communication:glands}{ग्रंथियाँ}, चिट्ठियाँ,
        पाने वाले --- एक सार-तालिका या अनुच्छेद में गिनाओ।
  \item अध्याय की सीख एक वाक्य में देकर बात ख़त्म करो: किस तरह का
        तंत्र किसी मान को बिना किसी चलाने वाले के, जीवन भर एक तंग गली
        में रखे रहता है?
\end{enumerate}
\end{problem}
